% --- Chapter 1: Introduction ---
\chapter{Introduction}
\label{chap:intro}

\section{Motivation}
\label{sec:motivation_intro}
Keeping our homes and workplaces clean is important for a healthy and pleasant life. But cleaning floors by hand often takes a lot of time, is boring, and can be hard physical work. In today's busy world, people have many things to do. Finding enough time and energy for regular cleaning can be tough, especially in big homes or busy business places. New technologies, especially in robotics, offer a good solution. Cleaning robots, which were once just an idea, are now helpful tools. They make cleaning easier, do a consistent job, and most importantly, save people's time. This saved time can be used for other important tasks, learning new things, or relaxing. The development of this technology has not only changed how we do chores but has also created new business opportunities and markets. For example, companies like iRobot, with its famous Roomba, were pioneers. They showed that many people wanted these new cleaning robots. After that, big home appliance companies like Samsung, LG, Dyson, and SharkNinja, and robot companies like Ecovacs and Roborock, also started making their own cleaning robots. They saw it as an important and growing part of the smart home market. This shows how automated cleaning has made a big impact, turning a new idea into a large global business.
This project looks at cleaning in Egyptian homes. In Egypt, many floors are made of tile or marble. Also, some areas can have a lot of dust. These conditions make automated cleaning robots very useful. Many advanced cleaning robots are sold worldwide, but they can be very expensive for some people. Because of this, and because people might want robots that fit their specific needs, there's a chance to create simpler, effective robots that can be made or put together locally. So, the main reason for this bachelor thesis is to design and build a smart cleaning robot that works well, is practical, and doesn't cost too much. The goal is to make a robot that greatly cuts down on the handwork needed for common floor cleaning, making homes more efficient and life easier. This project tries to find a good balance between useful features and cost, offering a helpful step in home robotics.

\section{Aim of the project}
\label{sec:aim_intro}
The primary technical aim of this work is to design and implement a functional prototype of a smart cleaning robot. This robot is intended to automate common floor cleaning tasks, thereby reducing manual effort and enhancing convenience within a typical household environment. The project focuses on delivering a practical and cost-effective solution that balances essential functionalities with simplicity of operation. The specific tasks to achieve this aim involve the implementation of the following core robot functions:
\begin{itemize}
    \item \textbf{Automated Vacuum Cleaning:} The robot will be equipped with a vacuuming system capable of collecting common household dust and light debris from floor surfaces. This function is central to its cleaning purpose.
    \item \textbf{Enhanced Autonomous Navigation:} The robot will implement a hybrid feed-forward and \hyperref[term:pid]{PID} control system using wheel encoders. The aim is to significantly improve movement precision, enabling the robot to more reliably execute a systematic room coverage pattern based on its sensor feedback. The system will also include autonomous stall detection and a recovery sequence.
    \item \textbf{Remote Control Operation:} The robot will feature a \hyperref[term:ble]{BLE} remote-control interface, enabling users to manually guide its movements. This is enhanced by \hyperref[term:pid]{PID} control to maintain a straight path during forward motion and an automatic obstacle-stop safety feature.
    \item \textbf{User Feedback System:} The robot will incorporate basic visual indicators (e.g., an \hyperref[term:led]{LED}) to provide feedback to the user about its current operational status.
\end{itemize}
The successful implementation and qualitative verification of these functions will demonstrate the achievement of the project's technical aim. This project aims to deliver a foundational smart cleaning robot prototype, with the understanding that highly advanced features like sophisticated mapping (\hyperref[term:slam]{SLAM}) are outside the current scope.

\section{Thesis Structure}
\label{sec:structure_intro_section}
This thesis is organized into five chapters. Each chapter discusses a different part of the project, providing a clear and logical progression from initial concepts to final outcomes and future possibilities: % Updated chapter count

\begin{itemize}
    \item \textbf{\hyperref[chap:intro]{Chapter 1: Introduction (This Chapter)}} This chapter provides the background for the project. It explains the motivation for developing a smart cleaning robot, outlines the project's objectives, and details the organization of this document.
    \item \textbf{\hyperref[chap:background]{Chapter 2: Theoretical Background and Literature Review}} This chapter delves into the technologies, design principles, and prior research relevant to autonomous and semi-autonomous cleaning robots. It examines various non-commercial and commercial robot designs, focusing on their navigation strategies, sensing capabilities, cleaning mechanisms, and control systems. It also reviews specific path planning techniques like the boustrophedon pattern.
    \item \textbf{\hyperref[chap:mechanical_design_assembly]{Chapter 3: Mechanical Design and Assembly}} This chapter focuses on the physical construction of the robot. It details the selection of mechanical components, the SolidWorks-based design of the chassis and custom parts, the selection of materials, the mechanical assembly process, and the specifics of the vacuum system's mechanical integration.
    \item \textbf{\hyperref[chap:electrical_control_performance]{Chapter 4: Electrical and Control System Design and Testing}} This chapter describes the electronic ``brains and nerves'' of the robot. It covers the selection of electrical components, the overall system architecture centered on the Raspberry Pi Pico W, the detailed electrical design and integration including the custom \hyperref[term:pcb]{PCB}, and the embedded software development for control. Finally, it presents a qualitative performance evaluation of the assembled prototype.
    \item \textbf{\hyperref[chap:conclusion_future_work_new]{Chapter 5: Conclusion and Recommendations for Future Work}} This chapter summarizes the project's main achievements and findings. It reflects on the extent to which the initial goals were met, discusses the identified limitations, and proposes recommendations for future research and potential improvements to enhance the robot's functionality and performance.
\end{itemize}

This structured approach aims to provide a comprehensive account of this project from its inception to its current state.